\documentclass[11pt]{article}
\usepackage{explorer}
\newcommand{\ExplorerLanguage}{FR}
\newcommand{\ExplorerNumber}{07}
\newcommand{\ExplorerKind}{Problème de synthèse}
\begin{document}
\ExplorerTitle{Première--Terminale · probabilités}{Inverser une probabilité conditionnelle}
\ExplorerSection{Données}
On choisit A avec probabilité \(\frac{1}{3}\) et B avec probabilité \(\frac{2}{3}\). L’urne A contient 3 rouges et 1 bleue ; B contient 1 rouge et 3 bleues. Pour une seule boule rouge observée, la probabilité d’avoir choisi A est \(\frac{3}{5}\).\par
\vspace{5pt}
\ExplorerSection{Chercher, justifier, conclure}
On reprend les deux urnes de l’exemple. Après avoir choisi une urne une seule fois, on effectue deux tirages avec remise dans cette même urne. Calculer la probabilité de deux rouges, puis celle d’avoir choisi A sachant que les deux boules sont rouges. Comparer avec l’observation d’une seule rouge et justifier les produits utilisés.\par
\vspace{5pt}
\end{document}
